\begin{frame}{Por que o oscilador harmônico?}
  \begin{columns}[T,onlytextwidth]
    \begin{column}{0.54\textwidth}
      \begin{itemize}
        \item Perto de um mínimo local com \alert{$V''(x_0)>0$},
              o potencial é aproximadamente parabólico:
      \end{itemize}
      \begin{equacao}
        \[
          V(x) \approx V(x_0)
          + \tfrac12 V''(x_0)\,(x-x_0)^2
        \]
      \end{equacao}
      \begin{itemize}
        \item Por isso ele descreve \alert{moléculas}, \alert{sólidos},
              \alert{campos} e \alert{circuitos}.
        \item É um dos poucos modelos quânticos com solução
              \alert{analítica exata}.
      \end{itemize}
    \end{column}
    \begin{column}{0.42\textwidth}
      \begin{destaque}{Onde ele aparece}
        \begin{itemize}
          \item Vibrações moleculares (espectroscopia IV)
          \item Fônons em cristais
          \item Íons aprisionados e qubits
          \item Ressonadores nanomecânicos
        \end{itemize}
      \end{destaque}
    \end{column}
  \end{columns}
\end{frame}
